\section*{Capítulo \ref{ch:g11:magnetic-fields} --- Magnetismo e campos magnéticos}

\begin{solution}{exo:g11:magnetic-fields:1}
Cada pedaço é um \omterm{def:g11:magnetic-fields:magnet}{ímã} completo, com um polo norte e um polo sul: um polo
novo aparece em cada face cortada. As faces recém-cortadas são um N e um
S, de modo que se atraem --- os pedaços tentam se remontar. Esquema
algum isola um polo: todo corte cria um par.
\end{solution}

\begin{solution}{exo:g11:magnetic-fields:2}
(a) Falso: o \omterm{def:g11:electric-gravitational-fields:field}{campo} tem um único sentido em cada ponto. (b) Verdadeiro.
(c) Verdadeiro: cada linha se fecha através do corpo do \omterm{def:g11:magnetic-fields:magnet}{ímã}. (d) Falso: a
agulha se acomoda \emph{ao longo} da linha, tangente a ela.
\end{solution}

\begin{solution}{exo:g11:magnetic-fields:3}
$B = \num{2e-7} \times 10 / 0.020 = \qty{1.0e-4}{T}$ --- o dobro do \omterm{def:g11:electric-gravitational-fields:field}{campo}
da Terra.
\end{solution}

\begin{solution}{exo:g11:magnetic-fields:4}
$n = 800/0.40 = \qty{2000}{m^{-1}}$;
$B = \num{1.26e-6} \times 2000 \times 1.5 \approx \qty{3.8}{mT}$. Como
$B \propto I$, dobrar a corrente dá \qty{7.5}{mT}.
\end{solution}

\begin{solution}{exo:g11:magnetic-fields:5}
$F = ILB = 5.0 \times 0.10 \times 0.20 = \qty{0.10}{N}$ --- o \omterm{def:g10:universal-gravitation:weight}{peso} de
$m = F/g = 0.10/9.81 \approx \qty{10}{g}$.
\end{solution}

\begin{solution}{exo:g11:magnetic-fields:6}
\emph{1.} A rota dela aponta $\ang{8}$ a oeste do norte verdadeiro:
$2000 \times \sin\ang{8} \approx \qty{2.8e2}{m}$ para oeste.

\emph{2.} Seguir $\ang{8}$ a \emph{leste} do norte da agulha.
\end{solution}

\begin{solution}{exo:g11:magnetic-fields:7}
$d = \mu_0 I / (2\pi B) = \num{2e-7} \times 20 / \num{5.0e-5}
= \qty{8.0}{cm}$. A menos de um decímetro de um cabo desses, a bússola lê
o fio, e não o planeta --- mantenha-a longe de condutores energizados.
\end{solution}

\begin{solution}{exo:g11:magnetic-fields:8}
$n = B/(\mu_0 I) = 0.010 / (\num{1.26e-6} \times 4.0)
\approx \qty{2.0e3}{m^{-1}}$; em \qty{25}{cm},
$N = 1989 \times 0.25 \approx 500$ espiras.
\end{solution}

\begin{solution}{exo:g11:magnetic-fields:9}
$F = 10 \times 0.12 \times 0.50 = \qty{0.60}{N}$;
$a = F/m = 0.60/0.020 = \qty{30}{m/s^2}$ --- cerca de $3g$: a barra salta
mesmo.
\end{solution}

\begin{solution}{exo:g11:magnetic-fields:10}
Equilíbrio: $ILB = mg$, logo
$I = \num{5.0e-3} \times 9.81 / (0.20 \times 0.10) \approx
\qty{2.5}{A}$. A \omterm{def:g10:inertia:force}{força} tem de apontar para cima; com $\vect B$
horizontal e perpendicular ao fio, a mão direita espalmada (palma para
cima) fixa o sentido necessário da corrente.
\end{solution}

\begin{solution}{exo:g11:magnetic-fields:11}
Mão direita espalmada: (a) para o alto da página; (b) para a direita;
(c) \omterm{def:g10:inertia:force}{força} nula --- o fio é paralelo ao \omterm{def:g11:electric-gravitational-fields:field}{campo}.
\end{solution}

\begin{solution}{exo:g11:magnetic-fields:12}
\emph{1.} $B = \num{2e-7} \times 5.0 / 0.030 \approx
\qty{33}{\micro T}$; as linhas circulam o fio, de modo que, logo abaixo
dele, o \omterm{def:g11:electric-gravitational-fields:field}{campo} é horizontal e perpendicular ao fio: leste--oeste.

\emph{2.} A agulha segue a soma vetorial: $\tan\theta = 33/20 = 1.67$,
logo $\theta \approx \ang{59}$ afastada do norte.
\end{solution}

\begin{solution}{exo:g11:magnetic-fields:13}
\emph{1.} $I = U/R = 12/2.0 = \qty{6.0}{A}$;
$n = 400/0.20 = \qty{2000}{m^{-1}}$;
$B_0 = \num{1.26e-6} \times 2000 \times 6.0 \approx \qty{15}{mT}$; com o
núcleo, $100 B_0 \approx \qty{1.5}{T}$.

\emph{2.} $P = UI = 12 \times 6.0 = \qty{72}{W}$.

\emph{3.} Sem corrente, sem \omterm{def:g11:electric-gravitational-fields:field}{campo} --- e o ferro doce não permanece
imantado: a carga se solta na hora, exatamente como se projetou.
\end{solution}

\begin{solution}{exo:g11:magnetic-fields:14}
\emph{1.} $F = NILB = 50 \times 2.0 \times 0.050 \times 0.30 =
\qty{1.5}{N}$ sobre cada lado. A corrente corre em sentidos opostos nos
dois lados, de modo que as \omterm{def:g10:inertia:force}{forças} são opostas: um binário que faz a
espira girar.

\emph{2.} Quando fica paralelo a $\vect B$, um fio não sente \omterm{prop:g11:magnetic-fields:laplace}{força de
Laplace} alguma.

\emph{3.} Depois de meia volta o binário se inverteria e desfaria a
rotação; o \omterm{rem:g11:magnetic-fields:motor}{comutador} troca o sentido da corrente a cada meia volta, de
modo que a espira continua girando sempre para o mesmo lado.
\end{solution}

\begin{solution}{exo:g11:magnetic-fields:15}
\emph{1.} $d = \num{2e-7} \times 100 / 1 = \qty{2.0e-5}{m} =
\qty{20}{\micro m}$ --- bem \emph{dentro} do cobre de um milímetro de
espessura. Nenhum ponto acessível perto de um fio isolado chega a
\qty{1}{T}: os \omterm{def:g11:electric-gravitational-fields:field}{campos} fortes empilham milhares de espiras (um
\omterm{def:g11:magnetic-fields:solenoid}{solenoide}), mais ferro ou supercondutores.

\emph{2.} $\num{e8} / \num{5e-5} = \num{2e12}$: cerca de doze potências
de dez.
\end{solution}

\begin{solution}{pb:g11:magnetic-fields:1}
\textbf{1.} A agulha é, ela mesma, um pequeno \omterm{def:g11:magnetic-fields:magnet}{ímã}; o \omterm{def:g11:electric-gravitational-fields:field}{campo} da Terra a
alinha, com a ponta norte ao longo da componente horizontal de $\vect B$,
ou seja, em direção ao norte magnético.

\textbf{2.} $B = \sqrt{20^2 + 44^2} = \sqrt{2336} \approx
\qty{48}{\micro T}$, inclinado $\tan^{-1}(44/20) \approx \ang{66}$
abaixo da horizontal.

\textbf{3.} A rota corre $\ang{6}$ a oeste do norte verdadeiro:
$30 \times \sin\ang{6} \approx \qty{3.1}{km}$ para oeste.

\textbf{4.} Os erros se somam: $\ang{10}$, logo
$30 \times \sin\ang{10} \approx \qty{5.2}{km}$ --- cinco quilômetros
para dez graus.

\textbf{5.} A componente horizontal: perto do polo o \omterm{def:g11:electric-gravitational-fields:field}{campo} é quase
vertical, a parte horizontal encolhe até quase zero e a agulha não tem
mais com o que se alinhar.

\textbf{6.} $n = 600/0.30 = \qty{2000}{m^{-1}}$.

\textbf{7.} $I = U/R = 12/1.5 = \qty{8.0}{A}$.

\textbf{8.} $B_0 = \mu_0 n I = \num{1.26e-6} \times 2000 \times 8.0
\approx \qty{20}{mT}$.

\textbf{9.} $50 \times \qty{20}{mT} = \qty{1.0}{T}$ --- o dobro do \omterm{def:g11:electric-gravitational-fields:field}{campo}
no polo de um \omterm{def:g11:magnetic-fields:magnet}{ímã} de neodímio, sobre uma face bem maior.

\textbf{10.} $P = UI = 12 \times 8.0 = \qty{96}{W}$; em
\qty{10}{\minute}: $E = 96 \times 600 \approx \qty{58}{kJ}$.

\textbf{11.} Porque ele solta: abra a chave, o \omterm{def:g11:electric-gravitational-fields:field}{campo} morre e o carro cai
exatamente onde se quer. Um \omterm{def:g11:magnetic-fields:magnet}{ímã} permanente agarra para sempre.

\textbf{12.} $F = NILB = 100 \times 3.0 \times 0.040 \times 0.25 =
\qty{3.0}{N}$.

\textbf{13.} Um binário: aplicadas em lados opostos do eixo, as duas
\omterm{def:g10:inertia:force}{forças} opostas giram a espira no mesmo sentido. \omterm{def:g10:inertia:force}{Forças} iguais e
\emph{paralelas} só empurrariam a espira de lado, sem fazê-la girar.

\textbf{14.} Nenhuma: um fio paralelo ao \omterm{def:g11:electric-gravitational-fields:field}{campo} não sente \omterm{prop:g11:magnetic-fields:laplace}{força de
Laplace}.

\textbf{15.} Ele inverte a corrente na espira a cada meia volta,
exatamente quando o binário mudaria de sinal --- de modo que o torque
sempre impulsiona a mesma rotação.

\textbf{16.} Dobrar $N$, ou dobrar $I$, ou dobrar $B$: o par de \omterm{def:g10:inertia:force}{forças}
$F = NILB$ é proporcional a cada um deles.

\textbf{17.} $1.5 / \num{5.0e-5} = \num{3.0e4}$: trinta mil Terras.

\textbf{18.} $n = B/(\mu_0 I) = 1.5 / (\num{1.26e-6} \times 500)
\approx \qty{2.4e3}{m^{-1}}$ --- e cada uma dessas espiras conduz
\qty{500}{A}.

\textbf{19.} $P = RI^2 = 0.20 \times 500^2 = \qty{50}{kW}$ --- o
aquecimento de um bairro inteiro. O enrolamento real é supercondutor:
\omterm{def:g11:circuits-and-power:resistance}{resistência} nula, dissipação nula, desde que mantido a poucos graus
acima do zero absoluto.

\textbf{20.} Bússola $\num{5e-5}$, guindaste $1$, ressonância $1.5$: o
navio fica cinco degraus abaixo das duas máquinas, que dividem uma mesma
década --- entre o $\num{e-10}$ interestelar e o $\num{e8}$ da estrela de
nêutrons. Nada mudou além dos \omterm{def:g11:circuits-and-power:current}{ampères} (e das espiras, do ferro e do
supercondutor que os multiplicam): um \omterm{def:g11:electric-gravitational-fields:field}{campo}, três trabalhos.
\end{solution}
